\documentclass[12pt,a4paper]{article}

% ─── Packages ───
\usepackage[utf8]{inputenc}
\usepackage[T1]{fontenc}
\usepackage[french]{babel}
\usepackage{geometry}
\geometry{margin=2.5cm}
\usepackage{graphicx}
\usepackage{listings}
\usepackage{xcolor}
\usepackage{booktabs}
\usepackage{hyperref}
\usepackage{amsmath}
\usepackage{tikz}
\usetikzlibrary{arrows.meta, positioning}

% ─── Style des listings C ───
\definecolor{codegreen}{rgb}{0,0.5,0}
\definecolor{codegray}{rgb}{0.5,0.5,0.5}
\definecolor{codepurple}{rgb}{0.58,0,0.82}
\definecolor{backcolour}{rgb}{0.97,0.97,0.97}

\lstdefinestyle{cstyle}{
  backgroundcolor=\color{backcolour},
  commentstyle=\color{codegreen},
  keywordstyle=\color{blue}\bfseries,
  numberstyle=\tiny\color{codegray},
  stringstyle=\color{codepurple},
  basicstyle=\ttfamily\footnotesize,
  breakatwhitespace=false,
  breaklines=true,
  captionpos=b,
  keepspaces=true,
  numbers=left,
  numbersep=5pt,
  showspaces=false,
  showstringspaces=false,
  showtabs=false,
  tabsize=2,
  language=C,
  morekeywords={pragma, omp, parallel, single, task, taskwait, taskgroup,
                shared, firstprivate, depend, final, num_threads, if},
  frame=single,
  rulecolor=\color{gray},
}
\lstset{style=cstyle}

\lstdefinestyle{bashstyle}{
  backgroundcolor=\color{backcolour},
  basicstyle=\ttfamily\footnotesize,
  breaklines=true,
  frame=single,
  rulecolor=\color{gray},
  language=bash,
}

% ─── Titre ───
\title{\textbf{TP4 --- Programmation OpenMP : Tâches}}
\author{}
\date{}

\begin{document}

\maketitle
\tableofcontents
\newpage

% ============================================================================
\section{Partie I --- Parcours d'un arbre (Depth First Search)}
% ============================================================================

\textbf{Fichier} : \texttt{Tree-traversal/tree-traverse.c}

Le programme construit un arbre binaire de recherche (BST) aléatoire et calcule
sa profondeur via un algorithme DFS (\textit{Depth First Search}). Toutes les
versions sont compilées et exécutées dans un même binaire pour faciliter la
comparaison.

% ----------------------------------------------------------------------------
\subsection{Q.1 --- Version séquentielle}
% ----------------------------------------------------------------------------

\textbf{Objectif} : Implémenter \texttt{inorderTraverse\_Seq(struct node *r)}
qui calcule la profondeur d'un arbre en séquentiel.

\textbf{Principe} : On parcourt récursivement les sous-arbres gauche et droit.
Chaque appel retourne la profondeur du sous-arbre. On renvoie
$1 + \max(\text{profondeur\_gauche}, \text{profondeur\_droite})$ pour compter le
niveau courant.

\begin{lstlisting}[caption={Version séquentielle --- Q.1}]
int inorderTraverse_Seq(struct node *r)
{
  int depthLeft = 0;
  int depthRight = 0;

  if (r->right != NULL)
  {
    depthRight = inorderTraverse_Seq(r->right);
  }
  if (r->left != NULL)
  {
    depthLeft = inorderTraverse_Seq(r->left);
  }

  /* +1 pour compter le niveau courant */
  return (depthLeft >= depthRight) ? depthLeft + 1 : depthRight + 1;
}
\end{lstlisting}

\textbf{Explications} :
\begin{itemize}
  \item \textbf{Cas de base} : si un n\oe ud est une feuille (pas d'enfants),
    \texttt{depthLeft} et \texttt{depthRight} restent à 0, et la fonction
    retourne 1 (la feuille elle-même).
  \item \textbf{Cas récursif} : on descend dans chaque sous-arbre non-nul, puis
    on prend le maximum des deux profondeurs $+ 1$.
\end{itemize}

% ----------------------------------------------------------------------------
\subsection{Q.2 --- Parallélisation avec tâches OpenMP}
% ----------------------------------------------------------------------------

\textbf{Objectif} : Paralléliser le parcours DFS avec \texttt{\#pragma omp task}.

\textbf{Principe} : Chaque appel récursif est encapsulé dans une tâche OpenMP.
Les deux sous-arbres (gauche et droit) peuvent être explorés en parallèle par
des threads différents. Un \texttt{\#pragma omp taskwait} synchronise les
résultats avant de calculer le maximum.

\begin{lstlisting}[caption={Version parallèle avec tâches --- Q.2}]
int inorderTraverse_Task(struct node *r)
{
  int depthLeft = 0;
  int depthRight = 0;

  if (r->right != NULL)
  {
    #pragma omp task shared(depthRight)
    {
      depthRight = inorderTraverse_Task(r->right);
    }
  }
  if (r->left != NULL)
  {
    #pragma omp task shared(depthLeft)
    {
      depthLeft = inorderTraverse_Task(r->left);
    }
  }

  #pragma omp taskwait

  return (depthLeft >= depthRight) ? depthLeft + 1 : depthRight + 1;
}
\end{lstlisting}

\textbf{Points clés de l'implémentation} :
\begin{itemize}
  \item \textbf{\texttt{shared(depthRight)} / \texttt{shared(depthLeft)}} : les
    variables de résultat doivent être partagées pour que la tâche puisse écrire
    le résultat et que le thread parent puisse le lire.
  \item \textbf{\texttt{\#pragma omp taskwait}} : point de synchronisation
    indispensable --- on doit attendre que les deux tâches enfants aient terminé
    avant de comparer les profondeurs.
  \item \textbf{Lancement} : une région \texttt{parallel} avec \texttt{single}
    est utilisée dans le \texttt{main} pour que seul un thread crée la première
    tâche, puis les tâches se répartissent naturellement entre les threads.
\end{itemize}

\begin{lstlisting}[caption={Lancement dans le main}]
#pragma omp parallel num_threads(numThreads)
{
  #pragma omp single
  {
    depth = inorderTraverse_Task(root);
  }
}
\end{lstlisting}

% ----------------------------------------------------------------------------
\subsection{Q.3 --- Limitation de la profondeur de génération de tâches}
% ----------------------------------------------------------------------------

\textbf{Problème} : La version Q.2 crée une tâche par n\oe ud de l'arbre. Pour
1~million de n\oe uds, cela représente $\sim$2~millions de tâches, ce qui génère
un surcoût considérable (création, ordonnancement, synchronisation).

\textbf{Solution} : Limiter la création de tâches aux \texttt{MAX\_DEPTH}
premiers niveaux de récursion ($2^5 = 32$ tâches maximum). Au-delà, le parcours
redevient séquentiel.

\subsubsection{Méthode 1 --- Clause \texttt{if}}

\begin{lstlisting}[caption={Limite de profondeur avec clause \texttt{if} --- Q.3a}]
#define MAX_DEPTH 5

int inorderTraverse_Task_DepthLimit_If(struct node *r, int currentDepth)
{
  int depthLeft = 0;
  int depthRight = 0;

  if (r->right != NULL)
  {
    #pragma omp task shared(depthRight) if(currentDepth < MAX_DEPTH)
    {
      depthRight = inorderTraverse_Task_DepthLimit_If(r->right,
                                                       currentDepth + 1);
    }
  }
  if (r->left != NULL)
  {
    #pragma omp task shared(depthLeft) if(currentDepth < MAX_DEPTH)
    {
      depthLeft = inorderTraverse_Task_DepthLimit_If(r->left,
                                                      currentDepth + 1);
    }
  }

  #pragma omp taskwait

  return (depthLeft >= depthRight) ? depthLeft + 1 : depthRight + 1;
}
\end{lstlisting}

\textbf{Fonctionnement de \texttt{if(condition)}} :
\begin{itemize}
  \item Si la condition est \textbf{vraie} : une tâche est créée normalement
    (peut être exécutée par un autre thread).
  \item Si la condition est \textbf{fausse} (\texttt{if(0)}) :
    \textbf{aucune tâche n'est créée}. Le code du bloc est exécuté immédiatement
    par le thread courant, comme un appel de fonction classique.
\end{itemize}

\subsubsection{Méthode 2 --- Clause \texttt{final}}

\begin{lstlisting}[caption={Limite de profondeur avec clause \texttt{final} --- Q.3b}]
int inorderTraverse_Task_DepthLimit_Final(struct node *r, int currentDepth)
{
  int depthLeft = 0;
  int depthRight = 0;

  if (r->right != NULL)
  {
    #pragma omp task shared(depthRight) final(currentDepth >= MAX_DEPTH)
    {
      depthRight = inorderTraverse_Task_DepthLimit_Final(r->right,
                                                          currentDepth + 1);
    }
  }
  if (r->left != NULL)
  {
    #pragma omp task shared(depthLeft) final(currentDepth >= MAX_DEPTH)
    {
      depthLeft = inorderTraverse_Task_DepthLimit_Final(r->left,
                                                         currentDepth + 1);
    }
  }

  #pragma omp taskwait

  return (depthLeft >= depthRight) ? depthLeft + 1 : depthRight + 1;
}
\end{lstlisting}

\textbf{Fonctionnement de \texttt{final(condition)}} :
\begin{itemize}
  \item Si la condition est \textbf{vraie} : la tâche est créée mais marquée
    \textbf{finale}. Elle et toutes ses tâches descendantes seront exécutées
    \textbf{immédiatement} (séquentiellement).
  \item \textbf{Différence avec \texttt{if}} : avec \texttt{final},
    \textbf{un objet tâche est toujours créé} (même s'il s'exécute
    immédiatement), ce qui engendre un léger surcoût par rapport à
    \texttt{if(0)} qui ne crée aucune tâche.
\end{itemize}

% ----------------------------------------------------------------------------
\subsection{Q.4 --- Analyse des performances}
% ----------------------------------------------------------------------------

\subsubsection{Résultats expérimentaux}

\textbf{Configuration} : 1\,000\,000 n\oe uds, 4 threads.

\begin{table}[h!]
\centering
\begin{tabular}{lcc}
\toprule
\textbf{Version} & \textbf{Temps (sec)} & \textbf{Speedup vs séq.} \\
\midrule
Q.1 --- Séquentielle          & \textbf{0.024} & 1.0$\times$              \\
Q.2 --- Tasks (sans limite)   & 0.295          & 0.08$\times$ (12$\times$ plus lent) \\
Q.3a --- \texttt{if} (limit=5)    & \textbf{0.059} & 0.41$\times$             \\
Q.3b --- \texttt{final} (limit=5) & 0.144          & 0.17$\times$             \\
Q.5 --- \texttt{taskgroup}        & 0.312          & 0.08$\times$             \\
\bottomrule
\end{tabular}
\caption{Performances du parcours d'arbre --- 1M n\oe uds, 4 threads}
\end{table}

\subsubsection{Analyse}

\textbf{1. La version tasks sans limite est bien plus lente que la version
séquentielle.}

Le surcoût de création/ordonnancement de $\sim$2~millions de tâches est
\textbf{largement supérieur} au travail utile de chaque tâche (une simple
comparaison). Le grain de calcul par tâche est extrêmement fin : chaque tâche ne
fait qu'une addition et une comparaison. Le rapport travail/surcoût est
catastrophique.

\textbf{2. La clause \texttt{if} offre la meilleure performance parallèle.}

En limitant à 5 niveaux, on ne crée que $2^5 = 32$ tâches au maximum. Chaque
tâche parcourt ensuite un sous-arbre complet séquentiellement. Le surcoût est
minimal et on répartit efficacement le travail sur les 4 threads.

\textbf{3. La clause \texttt{final} est moins performante que \texttt{if}.}

Avec \texttt{final}, des objets tâches sont toujours créés à chaque n\oe ud
(même si elles s'exécutent immédiatement). Cela génère un surcoût de gestion
mémoire absent avec \texttt{if(0)}.

\textbf{4. Le séquentiel reste le plus rapide ici.}

Pour ce type de problème à grain très fin (le travail par n\oe ud est trivial),
le parallélisme de tâches n'apporte pas de gain. Le parcours d'arbre est
principalement limité par les accès mémoire (pointeurs) et le surcoût de
synchronisation dépasse le gain potentiel. Un arbre BST aléatoire est aussi
naturellement déséquilibré, ce qui limite la répartition de charge.

% ----------------------------------------------------------------------------
\subsection{Q.5 --- \texttt{taskgroup} vs \texttt{taskwait}}
% ----------------------------------------------------------------------------

\begin{lstlisting}[caption={Version avec \texttt{taskgroup} --- Q.5}]
int inorderTraverse_Taskgroup(struct node *r)
{
  int depthLeft = 0;
  int depthRight = 0;

  #pragma omp taskgroup
  {
    if (r->right != NULL)
    {
      #pragma omp task shared(depthRight)
      {
        depthRight = inorderTraverse_Taskgroup(r->right);
      }
    }
    if (r->left != NULL)
    {
      #pragma omp task shared(depthLeft)
      {
        depthLeft = inorderTraverse_Taskgroup(r->left);
      }
    }
  } /* Fin du taskgroup : toutes les taches descendantes terminees */

  return (depthLeft >= depthRight) ? depthLeft + 1 : depthRight + 1;
}
\end{lstlisting}

\subsubsection{Différence entre \texttt{taskwait} et \texttt{taskgroup}}

\begin{table}[h!]
\centering
\begin{tabular}{p{2.5cm}p{5.5cm}p{5.5cm}}
\toprule
 & \textbf{\texttt{taskwait}} & \textbf{\texttt{taskgroup}} \\
\midrule
\textbf{Portée}
  & Attend les tâches enfants \textbf{directes} uniquement (1 niveau)
  & Attend \textbf{toutes} les tâches descendantes créées dans le bloc
    (enfants, petits-enfants, etc.) \\
\midrule
\textbf{Syntaxe}
  & Directive simple (pas de bloc)
  & Directive structurée (avec un bloc \texttt{\{ \}}) \\
\midrule
\textbf{Granularité}
  & Fine : ne synchronise qu'un niveau
  & Large : synchronise tout le sous-arbre de tâches \\
\midrule
\textbf{Usage typique}
  & Quand chaque niveau gère sa propre synchronisation
  & Quand on veut garantir la complétion de tout un sous-graphe de tâches \\
\bottomrule
\end{tabular}
\caption{Comparaison \texttt{taskwait} vs \texttt{taskgroup}}
\end{table}

\textbf{Dans notre cas}, les deux fonctionnent correctement car chaque niveau de
récursion fait son propre \texttt{taskwait}/\texttt{taskgroup}. Cependant,
\texttt{taskgroup} est plus \textbf{robuste} : même si on oubliait un
\texttt{taskwait} à un niveau intermédiaire, le \texttt{taskgroup} au niveau
supérieur garantirait l'attente de toutes les tâches descendantes.

\textbf{En termes de performance}, \texttt{taskgroup} a un surcoût légèrement
supérieur car il doit traquer toutes les tâches descendantes (pas seulement les
enfants directs).

% ============================================================================
\newpage
\section{Partie II --- Stencil}
% ============================================================================

\textbf{Fichier de base} : \texttt{STENCIL/stencil\_seq.c}

Le programme crée un maillage 1D rempli de valeurs aléatoires et applique
itérativement :
\begin{enumerate}
  \item \textbf{Un flou} (blur) sur les extrémités gauche et droite du maillage
    (moyennage avec les voisins).
  \item \textbf{Une copie} de la partie centrale (pas de blur).
  \item \textbf{Une mise à l'échelle} (scale) sur tout le maillage.
\end{enumerate}
La boucle s'arrête quand le blur converge (toutes les différences $<$
\texttt{threshold}).

% ----------------------------------------------------------------------------
\subsection{Q.6 --- Compilation et exécution}
% ----------------------------------------------------------------------------

\begin{lstlisting}[style=bashstyle, caption={Compilation et exécution}]
$ make
gcc -O3  -o stencil_seq stencil_seq.c
gcc -O3 -fopenmp -o stencil_task stencil_task.c
gcc -O3 -fopenmp -o stencil_taskdep stencil_taskdep.c

$ ./stencil_seq 100 10
Running stencil on 100 element(s) with seed 10
DONE w/ 5 iter in 0.000001 s
    MESH 0.00 0.02 0.04 0.05 0.05 0.06 0.06 0.06 0.07 0.08
\end{lstlisting}

Le Makefile compile 3 exécutables :
\begin{itemize}
  \item \texttt{stencil\_seq} : version séquentielle (fournie)
  \item \texttt{stencil\_task} : version avec tâches + \texttt{taskwait}
  \item \texttt{stencil\_taskdep} : version avec tâches + \texttt{depend}
\end{itemize}

% ----------------------------------------------------------------------------
\subsection{Q.7 --- Parallélisation avec tâches OpenMP}
% ----------------------------------------------------------------------------

\subsubsection{Analyse des dépendances}

Avant de paralléliser, identifions les dépendances entre les 4 boucles :

\begin{lstlisting}[style=bashstyle, caption={Dépendances entre boucles}]
Iteration k :
  Loop 1 (blur gauche)  : lit mesh[] -> ecrit mesh_blur[1..N/10-2]
  Loop 2 (copie milieu) : lit mesh[] -> ecrit mesh_blur[N/10-1..0.9N-2]
  Loop 3 (blur droit)   : lit mesh[] -> ecrit mesh_blur[0.9N-1..N-2]
  Loop 4 (scale)        : lit mesh_blur[] -> ecrit mesh[]
\end{lstlisting}

\textbf{Constat} :
\begin{itemize}
  \item \textbf{Loops 1, 2, 3} travaillent sur des régions \textbf{disjointes}
    de \texttt{mesh\_blur} et lisent toutes \texttt{mesh} $\rightarrow$
    \textbf{parallèles entre elles}.
  \item \textbf{Loop 4} lit \texttt{mesh\_blur} et écrit \texttt{mesh}
    $\rightarrow$ \textbf{dépend des loops 1, 2, 3}.
  \item \textbf{Itération k+1} lit \texttt{mesh} $\rightarrow$ \textbf{dépend
    de loop 4 de l'itération k}.
\end{itemize}

\begin{center}
\begin{tikzpicture}[
  node distance=1.2cm and 2.5cm,
  box/.style={draw, rounded corners, minimum width=3cm, minimum height=0.7cm,
              fill=blue!10, font=\footnotesize},
  arr/.style={-{Stealth[length=2.5mm]}, thick}
]
  \node[box] (mesh) {\texttt{mesh[]}};
  \node[box, right=of mesh, yshift=1.2cm]  (l1) {Loop 1 (blur gauche)};
  \node[box, right=of mesh]                (l2) {Loop 2 (copie milieu)};
  \node[box, right=of mesh, yshift=-1.2cm] (l3) {Loop 3 (blur droit)};
  \node[box, right=3cm of l2]              (l4) {Loop 4 (scale)};
  \node[box, right=of l4]                  (mesh2) {\texttt{mesh[]}};

  \draw[arr] (mesh) -- (l1);
  \draw[arr] (mesh) -- (l2);
  \draw[arr] (mesh) -- (l3);
  \draw[arr] (l1) -- (l4);
  \draw[arr] (l2) -- (l4);
  \draw[arr] (l3) -- (l4);
  \draw[arr] (l4) -- (mesh2);
  \draw[arr, dashed] (mesh2.south) -- ++(0,-0.8) -| node[below, pos=0.25, font=\footnotesize]{itération $k+1$} (mesh.south);
\end{tikzpicture}
\end{center}

\subsubsection{Version 1 : Tâches + \texttt{taskwait} (\texttt{stencil\_task.c})}

\begin{lstlisting}[caption={Stencil avec tâches + \texttt{taskwait} --- Q.7a}]
#pragma omp parallel
{
  #pragma omp single
  {
    do
    {
      int end1 = 1 ;
      int end3 = 1 ;

      /* Phase 1 : Blur (3 taches independantes) */

      /* Blur Kernel (loop 1) - cote gauche */
      #pragma omp task shared(mesh, mesh_blur, end1) \
                       firstprivate(N, threshold)
      {
        int ii ;
        for ( ii = 1 ; ii < N/10-1 ; ii++ )
        {
          double diff ;
          mesh_blur[ii] = (mesh[ii-1] + mesh[ii] + mesh[ii+1]) / 3 ;
          diff = mesh_blur[ii] - mesh[ii] ;
          if ( diff < 0 ) diff = -diff ;
          if ( diff > threshold ) end1 = 0 ;
        }
      }

      /* Blur Kernel (loop 2) - copie milieu */
      #pragma omp task shared(mesh, mesh_blur) firstprivate(N)
      {
        int ii ;
        for ( ii = N/10-1 ; ii < (int)(N*0.9-1) ; ii++ )
        {
          mesh_blur[ii] = mesh[ii] ;
        }
      }

      /* Blur Kernel (loop 3) - cote droit */
      #pragma omp task shared(mesh, mesh_blur, end3) \
                       firstprivate(N, threshold)
      {
        int ii ;
        for ( ii = (int)(N*0.9-1) ; ii < N - 1 ; ii++ )
        {
          double diff ;
          mesh_blur[ii] = (mesh[ii-1] + mesh[ii] + mesh[ii+1]) / 3 ;
          diff = mesh_blur[ii] - mesh[ii] ;
          if ( diff < 0 ) diff = -diff ;
          if ( diff > threshold ) end3 = 0 ;
        }
      }

      /* Synchronisation : attendre la fin des 3 blurs */
      #pragma omp taskwait

      end = end1 && end3 ;

      /* Phase 2 : Scale Kernel */
      #pragma omp task shared(mesh, mesh_blur) firstprivate(N)
      {
        int ii ;
        for ( ii = 0 ; ii < N ; ii++ )
        {
          mesh[ii] = mesh_blur[ii] * 3.3 / 5.4 ;
        }
      }

      /* Synchronisation : attendre le scale */
      #pragma omp taskwait

      niter++ ;
    }
    while( !end ) ;
  } /* fin single */
} /* fin parallel */
\end{lstlisting}

\textbf{Points clés} :
\begin{itemize}
  \item \textbf{\texttt{end1}, \texttt{end3} séparés} : on utilise deux
    variables locales au lieu de \texttt{end} directement pour éviter une data
    race entre les tâches 1 et 3 qui écrivent toutes deux le flag de
    convergence.
  \item \textbf{1\textsuperscript{er} \texttt{taskwait}} : indispensable avant
    d'accéder à \texttt{end1}/\texttt{end3} (écrits par les tâches) et avant de
    lancer le scale (qui lit \texttt{mesh\_blur}).
  \item \textbf{2\textsuperscript{ème} \texttt{taskwait}} : indispensable avant
    l'itération suivante (qui lit \texttt{mesh}, modifié par le scale).
  \item \textbf{\texttt{firstprivate(N, threshold)}} : chaque tâche reçoit une
    copie de ces scalaires pour éviter tout problème d'accès concurrent.
\end{itemize}

\subsubsection{Version 2 : Tâches + \texttt{depend} (\texttt{stencil\_taskdep.c})}

\begin{lstlisting}[caption={Stencil avec tâches + \texttt{depend} --- Q.7b}]
/* Indices des 3 regions pour les dependances */
int r1 = 1 ;
int r2 = N / 10 - 1 ;
int r3 = (int)(N * 0.9 - 1) ;

#pragma omp parallel
{
  #pragma omp single
  {
    do
    {
      int end1 = 1 ;
      int end3 = 1 ;

      /* Blur gauche */
      #pragma omp task shared(mesh, mesh_blur, end1) \
        firstprivate(N, threshold) \
        depend(in: mesh[0]) depend(out: mesh_blur[r1])
      {
        int ii ;
        for ( ii = 1 ; ii < N/10-1 ; ii++ )
        {
          double diff ;
          mesh_blur[ii] = (mesh[ii-1]+mesh[ii]+mesh[ii+1]) / 3 ;
          diff = mesh_blur[ii] - mesh[ii] ;
          if ( diff < 0 ) diff = -diff ;
          if ( diff > threshold ) end1 = 0 ;
        }
      }

      /* Blur milieu */
      #pragma omp task shared(mesh, mesh_blur) \
        firstprivate(N) \
        depend(in: mesh[0]) depend(out: mesh_blur[r2])
      {
        int ii ;
        for ( ii = N/10-1 ; ii < (int)(N*0.9-1) ; ii++ )
          mesh_blur[ii] = mesh[ii] ;
      }

      /* Blur droit */
      #pragma omp task shared(mesh, mesh_blur, end3) \
        firstprivate(N, threshold) \
        depend(in: mesh[0]) depend(out: mesh_blur[r3])
      {
        int ii ;
        for ( ii = (int)(N*0.9-1) ; ii < N-1 ; ii++ )
        {
          double diff ;
          mesh_blur[ii] = (mesh[ii-1]+mesh[ii]+mesh[ii+1]) / 3 ;
          diff = mesh_blur[ii] - mesh[ii] ;
          if ( diff < 0 ) diff = -diff ;
          if ( diff > threshold ) end3 = 0 ;
        }
      }

      /* Scale */
      #pragma omp task shared(mesh, mesh_blur) firstprivate(N) \
        depend(in: mesh_blur[r1], mesh_blur[r2], mesh_blur[r3]) \
        depend(out: mesh[0])
      {
        int ii ;
        for ( ii = 0 ; ii < N ; ii++ )
          mesh[ii] = mesh_blur[ii] * 3.3 / 5.4 ;
      }

      /* taskwait necessaire pour lire end1/end3 */
      #pragma omp taskwait

      end = end1 && end3 ;
      niter++ ;
    }
    while( !end ) ;
  }
}
\end{lstlisting}

\textbf{Avantages de \texttt{depend} par rapport à \texttt{taskwait}} :

\begin{enumerate}
  \item \textbf{Ordonnancement plus fin} : les dépendances sont exprimées entre
    tâches individuelles, pas entre phases entières. Le runtime peut
    potentiellement chevaucher le scale de l'itération $k$ avec les blurs de
    l'itération $k+1$.

  \item \textbf{Graphe de dépendances explicite} :
    \begin{itemize}
      \item Les 3 blurs ont \texttt{depend(in: mesh[0])} $\rightarrow$ ils
        attendent que le scale précédent ait fini d'écrire \texttt{mesh}.
      \item Le scale a \texttt{depend(in: mesh\_blur[r1], mesh\_blur[r2],
        mesh\_blur[r3])} $\rightarrow$ il attend les 3 blurs.
      \item Pas besoin de \texttt{taskwait} pour cet ordonnancement : les
        \texttt{depend} suffisent.
    \end{itemize}

  \item \textbf{Un \texttt{taskwait} reste nécessaire} uniquement pour lire
    \texttt{end1}/\texttt{end3} (le thread \texttt{single} doit attendre pour
    tester la condition de convergence de la boucle \texttt{do\ldots while}).
\end{enumerate}

\subsubsection{Résultats de performance}

\textbf{Configuration} : 1\,000\,000 éléments, seed 42, 4 threads.

\begin{table}[h!]
\centering
\begin{tabular}{lcc}
\toprule
\textbf{Version} & \textbf{Temps (sec)} & \textbf{Itérations} \\
\midrule
Séquentielle             & \textbf{0.008} & 5 \\
Tasks + \texttt{taskwait} & 0.042          & 5 \\
Tasks + \texttt{depend}   & \textbf{0.022} & 5 \\
\bottomrule
\end{tabular}
\caption{Performances du stencil --- 1M éléments, 4 threads}
\end{table}

\textbf{Analyse} :
\begin{itemize}
  \item Les 3 versions produisent le \textbf{même résultat} (MESH identique),
    validant la correction.
  \item Le séquentiel est le plus rapide car le nombre d'itérations est faible
    (5) et le surcoût de création des threads/tâches n'est pas amorti.
  \item La version \texttt{depend} est \textbf{2$\times$ plus rapide} que la
    version \texttt{taskwait} car le runtime dispose de plus d'informations pour
    ordonnancer les tâches efficacement.
  \item Pour des maillages plus grands ou plus d'itérations, l'avantage du
    parallélisme serait plus marqué, car les boucles de blur et scale
    deviendraient le goulot d'étranglement.
\end{itemize}

\end{document}
